%==============================================================================
%  27-predictions.tex — Predictions: AI Usage & Organizational Change
%==============================================================================
\strategyPart{23 \textbar\ Predictions}{Predictions of AI Usage and Organizational Change}

This part consolidates what leading institutions predict about \textbf{how
organizations will use AI and how they will change} on the basis of
generative AI — not market size. The pattern across the research is
consistent: usage will keep rising, and the differentiator will be
organizational change, not technology.

\section{Predictions of AI usage in organizations}

\begin{itemize}
  \item \textbf{Scaling is the next wave.} 88\% of organizations already use AI
        in at least one function, but only 7\% report full scaling — the
        coming wave is enterprise-wide deployment and workflow redesign
        [McKinsey State of AI 2025].
  \item \textbf{Generative AI becomes mainstream.} 65\% of organizations now
        regularly use generative AI in at least one function, up from roughly
        a third a year earlier [McKinsey State of AI 2024].
  \item \textbf{Agentic AI reaches the enterprise software stack.} Gartner
        predicts that by 2028, \textbf{33\% of enterprise software
        applications} will include agentic AI (up from under 1\% in 2024),
        enabling \textbf{15\% of day-to-day work decisions} to be made
        autonomously [Gartner 2024].
  \item \textbf{Consolidation after the experiment.} Through 2025,
        \textbf{30\% of generative-AI projects} will be abandoned after proof
        of concept as organizations separate what scales from what does not
        [Gartner 2023].
  \item \textbf{Workforce access broadens.} Deloitte's enterprise research
        documents the move from experimentation to scaling, with wider
        workforce access to AI and rising customization of autonomous agents
        [Deloitte State of AI in the Enterprise 2026].
\end{itemize}

\section{Predictions of organizational change}

The World Economic Forum's transformation research — based on insights from
more than \textbf{450 executives} — predicts change across five areas:
customer experience, operations and supply chains, R\&D, strategic planning
and talent. It emphasizes human accountability, end-to-end operating-model
redesign, scalable talent systems, transparency and disciplined
experimentation [WEF--Accenture 2026].

McKinsey's \textit{State of Organizations} research predicts that the
organizations which capture value will be those that redesign work
end-to-end — distributing responsibilities between humans, agents and
automation rather than adding AI to unchanged workflows [McKinsey 2025;
McKinsey Symbiotic Enterprise].

\begin{center}
  \fitwidth{\begin{tabular}{>{\raggedright\arraybackslash}p{3.2cm}>{\raggedright\arraybackslash}p{4.8cm}>{\raggedright\arraybackslash}p{3.6cm}}
    \rowcolor{inNavy}
    \hcell{Institution} & \hcell{Prediction} & \hcell{Horizon} \\
    \textbf{Gartner} & 33\% of enterprise software includes agentic AI; 15\% of decisions autonomous & 2028 \\
    \textbf{Gartner} & 30\% of gen-AI projects abandoned after proof of concept & 2025 \\
    \textbf{McKinsey} & 88\% use AI; only 7\% fully scaled — scaling becomes the priority & Current \\
    \textbf{Deloitte} & From experimentation to scaling; broader workforce access & 2026 \\
    \textbf{WEF} & End-to-end operating-model redesign across five areas & 2026 \\
    \textbf{Stanford} & Continued adoption growth; 78\% of organizations use AI in some function & 2025 \\
    \textbf{Microsoft} & 75\% of knowledge workers use AI at work; digital debt persists & 2024 \\
  \end{tabular}}
\end{center}

\section{The prediction that matters}

Every institution converges on the same forecast: \textbf{the value gap is
closing for organizations that redesign work, and widening for those that do
not}. The organizations that will change fastest are those already treating
generative AI as a problem-solving capability embedded in redesigned
workflows — with governance, skills and measurement in place — rather than as
a tool portfolio. This is the basis of the problem-first strategy in
Part~22.

\takeaway{The value gap closes for organizations that redesign work and widens for those that do not.}
